%%
%% Test: Book Structure (main.tex pattern) — v1.1
%% Stage: main structure (stages 1-4)
%% 
%% Purpose: Verify the full book structure:
%%   - Front cover
%%   - Title page
%%   - Copyright page
%%   - Dedication page
%%   - Preface
%%   - ToC, LoF, LoT
%%   - Main matter (chapters)
%%   - Bibliography
%%   - Index
%%   - Back cover
%% 

\documentclass{matinbook}

\addbibresource{../references.bib}

\title{تست ساختار کتاب}
\author{تیم MatinBook}
\date{۱۴۰۵}

\renewcommand{\repository}{github.com/matinbook/test}

\begin{document}

%----------------------------------------
% FRONT COVER
%----------------------------------------
\makecover
    {تست ساختار کتاب}
    {بررسی ترتیب اجزای کتاب}
    {تیم MatinBook}
    {۱۴۰۵}

%----------------------------------------
% FRONT MATTER
%----------------------------------------
\frontmatter

\maketitle

% Copyright page
\thispagestyle{empty}
\vfill
\begin{center}
    {\large\textbf{حقوق نشر}}
    
    \vspace{1cm}
    تمامی حقوق این کتاب محفوظ است.
    
    \vspace{1cm}
    \textbf{شناسنامه کتاب}
    
    \vspace{0.5cm}
    \begin{tabular}{rl}
        عنوان: & تست ساختار کتاب \\
        نویسنده: & تیم MatinBook \\
        ناشر: & MatinBook \\
        نوبت چاپ: & اول \\
        تیراژ: & ۱۰۰۰ نسخه \\
        شابک: & ۹۷۸-XXX-XXX-XXX-X \\
        قیمت: & XXX,XXX ریال \\
    \end{tabular}
    
    \vspace{1cm}
    {\small نسخه \matinbookversion\ — \matinbookdate}
\end{center}
\clearpage

% Dedication
\thispagestyle{empty}
\vfill
\begin{flushright}
    {\Large\itshape
        تقدیم به...
    }
\end{flushright}
\clearpage

% Preface
\chapter*{پیشگفتار}
\addcontentsline{toc}{chapter}{پیشگفتار}

این یک پیشگفتار نمونه است.

\clearpage

\tableofcontents
\clearpage

\listoffigures
\clearpage

\listoftables
\clearpage

%----------------------------------------
% MAIN MATTER
%----------------------------------------
\mainmatter

\chapter{فصل اول}
\label{chap:first}

این فصل اول است.

\section{بخش اول}
\label{sec:first-section}

این بخش اول است.

\begin{theorem}[قضیه نمونه]
    برای هر عدد حقیقی $x$، داریم $x^2 \geq 0$.
\end{theorem}

\begin{proof}
    چون مربع هر عدد حقیقی نامنفی است، نتیجه حاصل می‌شود.
\end{proof}

\chapter{فصل دوم}
\label{chap:second}

این فصل دوم است.

%----------------------------------------
% BACK MATTER
%----------------------------------------
\backmatter

\printbibliography[title={منابع و مراجع}]

\printindex

%----------------------------------------
% BACK COVER
%----------------------------------------
\makebackcover{%
    این یک تست برای بررسی ساختار کامل کتاب است.
}

\end{document}
